\section*{附录习题}
\phantomsection
\addcontentsline{toc}{section}{附录习题}

\begin{enumerate}
\def\labelenumi{\arabic{enumi})}
\tightlist
\item
  设 \(\mathfrak{K}\) 是一个哈代域，使得对每个函数 \(f\in \mathfrak{K}\) （在 \(+\infty\) 的某个邻域上不恒为零）存在一个 \(\lambda >0\) 使
\end{enumerate}

\[
\frac {1}{e _ {m} (x ^ {\lambda})} \ll f (x) \ll e _ {m} (x ^ {\lambda})
\]

（m 是一个与 f 无关的整数）。

\begin{enumerate}
\def\labelenumi{\alph{enumi})}
\tightlist
\item
  设 \(u_{1}, u_{2}, \ldots, u_{p}\) 是 p 个形如 \(u_{k} = \log |z_{k}|\) 的函数，其中 \(z_{k} \in K\) 在 \(+\infty\) 的某个邻域上不为零。证明：对每个函数 g（在 \(+\infty\) 的某个邻域上不为零；属于哈代域 \(\mathfrak{K}(u_{1}, u_{2}, \ldots, u_{p})\) ，该域由把诸函数 \(u_{k}\) （ \(1 \leqslant k \leqslant p\) ）添加到 K 而得到），存在一个数 \(\mu > 0\) 使
\end{enumerate}

\[
\frac {1}{e _ {m} (x ^ {\mu})} \ll g (x) \ll e _ {m} (x ^ {\mu})
\]

（归结为 g 是系数在 K 中的 \(u_{k}\) 的多项式的情形，对 p 作归纳论证；然后对 p = 1，仿照 V, p.~248 引理 2 的做法对多项式 g 的次数作归纳论证）。

\begin{enumerate}
\def\labelenumi{\alph{enumi})}
\setcounter{enumi}{1}
\tightlist
\item
  设 \(u_{k}\) （ \(1 \leqslant k \leqslant p\) ）是 \(p\) 个形如 \(u_{k} = \exp(z_{k})\) 的函数，其中 \(z_{k} \in \mathcal{K}\) 。证明：对每个函数 \(g\) （属于哈代域 \(\mathcal{K}(u_{1}, u_{2}, \ldots, u_{p})\) ，在 \(+\infty\) 的某个邻域上不恒为零），存在一个数 \(\mu > 0\) 使
\end{enumerate}

\[
\frac {1}{e _ {m + 1} (x ^ {\mu})} \prec g (x) \prec e _ {m + 1} (x ^ {\mu})
\]

（类似的方法）。

\begin{enumerate}
\def\labelenumi{\alph{enumi})}
\setcounter{enumi}{2}
\tightlist
\item
  推出：若 \(f\) 是一个具有 \(n\) 项定义序列、且在 \(+\infty\) 的某个邻域上不恒为零的 (H) 函数，则存在一个数 \(\lambda > 0\) 使
\end{enumerate}

\[
\frac {1}{e _ {n} (x ^ {\lambda})} \ll f (x) \ll e _ {n} (x ^ {\lambda}).
\]

\begin{enumerate}
\def\labelenumi{\arabic{enumi})}
\setcounter{enumi}{1}
\tightlist
\item
  \begin{enumerate}
  \def\labelenumii{\alph{enumii})}
  \tightlist
  \item
    证明每个只有一项定义序列的 (H) 函数等价于形如 \(x^{p}(\log x)^{q}\) 或 \(x^{p}e^{g(x)}\) 之一的函数，其中 p 与 q 是有理整数，g 是 x 的多项式（习题 1 的方法）。
  \end{enumerate}
\end{enumerate}

\begin{enumerate}
\def\labelenumi{\alph{enumi})}
\setcounter{enumi}{1}
\tightlist
\item
  由 a) 推出：哈代域 \(\mathbf{R}(x,u_{1},\ldots,u_{p})\) （其中 \(u_{1},\ldots,u_{p}\) 是只有单项定义序列的 (H) 函数）中每个函数都等价于形如 \(x^{p}(\log x)^{q}e^{g(x)}\) 的函数，其中 p 与 q 是有理整数，g 是 x 的多项式。
\end{enumerate}

\begin{enumerate}
\def\labelenumi{\arabic{enumi})}
\setcounter{enumi}{2}
\item
  设 f 与 g 是两个 (H) 函数，使 f/g 相对于 \(l_{m}(x)\) 的阶为 0；证明：若 g 相对于 \(l_{m}(x)\) 的阶不为 0，则 \(f'/g' \sim f/g\) （比较 \(\log |f|\) 与 \(\log |g|\) ）。
\item
  设 \(\mathfrak{K}\) 是一个哈代域，使得对每个函数 \(f\in \mathfrak{K}\) （不等价于常数、且相对于 \(l_{m - 1}(x)\) 的阶为 0），存在一个常数 \(k\) 与一个有理整数 \(r\) ，使 \(f(x)\sim k(l_m(x))^r\) 。
\end{enumerate}

\begin{enumerate}
\def\labelenumi{\alph{enumi})}
\item
  设 z 是 K 中任一在 \(+\infty\) 的某个邻域上不恒为零的函数。证明：哈代域 \(\mathcal{K}(\log|z|)\) 中每个不等价于常数、且相对于 \(l_{m}(x)\) 的阶为 0 的函数 g，都等价于形如 \(k(l_{m+1}(x))^{r}\) 的函数（k 为常数，r 为有理整数）。（先考虑 g 是 \(\log|z|\) 的 p 次多项式、系数在 K 中的情形，并用习题 3 对 p 作归纳论证；若 g 是 \(\log|z|\) 的有理函数、系数在 K 中，则用习题 3 对分子的次数作归纳论证。）
\item
  证明 \(\mathfrak{K}(l_{m+1}(x))\) 中每个函数都等价于形如 \(f(x)(l_{m+1}(x))^r\) 的函数，其中 \(f \in \mathfrak{K}\) ， \(r\) 是有理整数。
\item
  设 z 是 K 中任一函数。证明：哈代域 \(\mathfrak{K}(e^{z})\) 中每个相对于 \(l_{m}(x)\) 的阶为 0 的函数 g 都等价于一个常数。（先考虑 \(g = u^{q} \mathrm{P}(u)\) 的情形，其中 \(u = e^{z}\) ，q 是有理整数， \(\mathrm{P}(u)\) 是 u 的 p 次多项式、系数在 K 中；然后用 a) 与习题 3 对 p 作归纳论证。再用习题 3 过渡到一般情形。）
\item
  把 a) 的结果推广到域 \(\mathfrak{K}(u_{1}, u_{2}, \ldots, u_{s})\) ，其中诸 \(u_{k}\) 形如 \(e^{z_{k}}\) 或 \(\log |z_{k}|\) ，诸 \(z_{k}\) 是 K 中在 \(+\infty\) 的某个邻域上不恒为零的函数（对 s 作归纳论证）。
\end{enumerate}

\begin{enumerate}
\def\labelenumi{\arabic{enumi})}
\setcounter{enumi}{4}
\tightlist
\item
  \begin{enumerate}
  \def\labelenumii{\alph{enumii})}
  \tightlist
  \item
    由习题 4 推出：若 f 是一个具有 n 项定义序列、不等价于常数、且相对于 \(l_{n-1}(x)\) 的阶为 0 的 (H) 函数，则存在一个常数 k 与一个有理整数 r，使 \(f(x) \sim k(l_n(x))^r\) 。
  \end{enumerate}
\end{enumerate}

\begin{enumerate}
\def\labelenumi{\alph{enumi})}
\setcounter{enumi}{1}
\tightlist
\item
  由习题 4 推出：若 \(f\) 是任一 (H) 函数，则存在一个整数 \(n\) ，使 \(n\) 阶对数判别法能够判定积分 \(\int_{a}^{+\infty} f(t) dt\) 收敛或为无穷。
\end{enumerate}

\begin{enumerate}
\def\labelenumi{\arabic{enumi})}
\setcounter{enumi}{5}
\tightlist
\item
  把函数 \(e_p\left((l_q(x))^{\mu}\right)\) 按整数 \(p\) 与 \(q\) 以及实数 \(\mu\) （设其 \(> 0\) 且 \(\neq 1\) ）的取值相互比较。
\end{enumerate}

7) 设 \(f\) 是一个具有 \(n\) 项定义序列的 (H) 函数。证明：若有 \(e_p((l_{q+1}(x))^{\beta}) \ll f(x) \ll e_p((l_q(x))^{\alpha})\) （相应地， \(e_p((l_q(x))^{\beta}) \ll f(x) \ll e_{p+1}((l_q(x))^{\alpha}))\) ），

则对每个 \(\beta > 0\) 与每个 \(\alpha\) （ \(0 < \alpha < 1\) ），必有 \(p + q + 1 < n\) （对 \(n\) 作归纳论证，使用与习题 1 (V, p.~263) 和 4 (V, p.~264) 类似的方法）。

\begin{enumerate}
\def\labelenumi{\arabic{enumi})}
\setcounter{enumi}{7}
\tightlist
\item
  \begin{enumerate}
  \def\labelenumii{\alph{enumii})}
  \tightlist
  \item
    设 \((f_{n})\) 是属于 \(\mathcal{H}(\mathfrak{F},\mathbf{R})\) 的递增连续函数序列，且对每个 n 有 \(f_{n} \ll f_{n+1}\) 。证明存在一个连续、递增、属于 \(\mathcal{H}(\mathfrak{F},\mathbf{R})\) 的函数 f，使对每个 n 有 \(f \gg f_{n}\) （归结为 \(f_{n} \leqslant f_{n+1}\) 的情形，并定义 f 使 \(f_{n}(x) \leqslant f(x) \leqslant f_{n+1}(x)\) （对 \(n \leqslant x \leqslant n+1\) ））。
  \end{enumerate}
\end{enumerate}

\begin{enumerate}
\def\labelenumi{\alph{enumi})}
\setcounter{enumi}{1}
\item
  设 \((f_n)\) 是属于 \(\mathcal{H}(\mathfrak{F},\mathbf{R})\) 的递增连续函数序列，且 \(1\ll f_{n + 1}\ll f_n\) 对每个 \(n\) 成立。证明存在一个函数 \(f\) ，连续、递增、属于 \(\mathcal{H}(\mathfrak{F},\mathbf{R})\) ，使 \(1\ll f\ll f_n\) 对每个 \(n\) 成立（归结为 \(f_{n + 1}\leqslant f_n\) 的情形，并证明可以定义一个递增的实数序列 \((x_{n})\) 与一个连续递增函数 \(f\) ，使 \(f_{n + 1}(x)\leqslant f(x)\leqslant f_n(x)\) （对 \(x_{n}\leqslant x\leqslant x_{n + 1}\) ））。
\item
  设 \((f_n), (g_n)\) 是两个属于 \(\mathcal{H}(\mathfrak{F}, \mathbf{R})\) 的连续递增函数序列，满足 \(f_n \ll f_{n+1}\) 、 \(g_m \gg g_{m+1}\) 且 \(f_n \ll g_m\) （对所有 \(m\) 与 \(n\) ）；证明存在一个函数 \(h\) ，连续、递增、属于 \(\mathcal{H}(\mathfrak{F}, \mathbf{R})\) ，使 \(f_n \ll h \ll g_m\) （对所有 \(m\) 与 \(n\) ）（类似的方法）。
\end{enumerate}

特别地，证明存在一个连续递减函数 \(f\) （属于 \(\mathcal{H}(\mathfrak{F},\mathbf{R})\) ），使任何对数判别法都不能判定积分 \(\int_{a}^{+\infty}f(t)dt\) 收敛或为无穷（“杜·布瓦–雷蒙定理”）。

\begin{enumerate}
\def\labelenumi{\arabic{enumi})}
\setcounter{enumi}{8}
\item
  证明级数 \(\sum_{n=1}^{\infty} \frac{e_n(x)}{e_n(n)}\) 在 \(\mathbf{R}\) 的每个紧子集上一致收敛，且其和 \(f(x)\) 满足： \(f(x) \gg e_n(x)\) 对每个整数 \(n\) 成立。
\item
  证明存在一个递增函数 \(f\) ，它有定义、连续且 \(> 0\) （对 \(x \geqslant 0\) ），使 \(f(2x) = 2^{f(x)}\) 对所有 \(x \geqslant 0\) 成立；由此推出 \(f(x) \gg e_n(x)\) 对每个整数 \(n\) 成立。
\item
  设 \(f\) 是一个递增函数，连续且 \(\geqslant 0\) （对 \(x \geqslant 0\) 有定义），且 \(f(x) \gg e_n(x)\) 对每个整数 \(n\) 成立；证明：若 \(g\) 是 \(f\) 的反函数，则 \(1 \ll g(x) \ll l_n(x)\) 对每个整数 \(n\) 成立。
\end{enumerate}

12) 对每个整数 \(n > 0\) ，设 \(n = \sum_{k=0}^{\infty} \varepsilon_k 2^k\) 是 \(n\) 的二进展开（ \(\varepsilon_k\) 是整数，除有限个指标外均为零， \(0 \leqslant \varepsilon_k \leqslant 1\) ）。令

\[
\alpha (n) = \sum_ {k = 0} ^ {\infty} \varepsilon_ {k}, \quad \mathrm{A} _ {1} (n) = \sum_ {j = 1} ^ {n} \alpha (j), \quad \mathrm{A} _ {2} (n) = \sum_ {j = 1} ^ {n} \alpha (\alpha (j)).
\]

\begin{enumerate}
\def\labelenumi{\alph{enumi})}
\tightlist
\item
  证明公式 \(\mathrm{A}_1(n) = \sum_{k=0}^{\infty}(k+2)\varepsilon_k 2^{k-1}\) 。由此推出公式
\end{enumerate}

\[
\mathrm{A} _ {1} (n) = \frac {n}{2} \frac {\log n}{\log 2} + o (n \log \log n)
\]

其中 \(n\) 无限增大（把表示 \(\mathrm{A}_1(n)\) 的和分解成两部分：第一个和中指标 \(k\) 从 0 变到 \(\mu(n)\) ，第二个和中从 \(\mu(n)\) 变到 \(n_1\) ，这里 \(n_1\) 是数 \(k\) （使 \(\varepsilon_k \neq 0\) ）中最大者，而 \(\mu(n)\) 取得适当；用 V, p.~239 的命题 6 估计第一个和的上界，然后估计第二个和与 \(n_1n / 2\) 之差。

\begin{enumerate}
\def\labelenumi{\alph{enumi})}
\setcounter{enumi}{1}
\tightlist
\item
  建立关系式
\end{enumerate}

\[
\sum_ {j = 0} ^ {2 ^ {m} - 1} f (\alpha (j)) = \sum_ {k = 0} ^ {m} \binom {m} {k} f (k)
\]

其中 \(f\) 是定义在 \(\mathbf{N}\) 上的任意函数（对给定的 \(k\) ，考察 \(j < 2^{m}\) 中使 \(\alpha(j) = k\) 的个数）。

\begin{enumerate}
\def\labelenumi{\alph{enumi})}
\setcounter{enumi}{2}
\tightlist
\item
  对 \(m = 2^r - 1\) ，证明 \(\alpha(m - j) + \alpha(j) = r\) 。由此并借助 \(b\) ) 推出
\end{enumerate}

\[
\mathrm{A} _ {2} (2 ^ {m}) = r 2 ^ {m - 1} \sim \frac {2 ^ {m} \log \log 2 ^ {m}}{2 \log 2}.
\]

\begin{enumerate}
\def\labelenumi{\alph{enumi})}
\setcounter{enumi}{3}
\tightlist
\item
  对 \(m = 2^{r}\) ，证明
\end{enumerate}

\[
\mathrm{A} _ {2} (2 ^ {m}) = 2 \mathrm{A} _ {2} (2 ^ {m - 1}) + 2 ^ {m - 1} - \sum_ {j = 0} ^ {m - 1} \binom {m - 1} {j} \Delta (j + 1)
\]

其中我们令 \(\alpha(n+1)-\alpha(n)=1-\Delta(n+1)\) （对所有 \(n\) ；利用 \(b\) ）以及关系式 \(\alpha(2^{k}+r)=\alpha(r)+1\) （对 \(r<2^{k}\) ））。证明 \(\sum_{j=0}^{m-1}\binom{m-1}{j}\Delta(j+1)\leqslant r^{m-2}\) （注意 \(\Delta(j+1)=0\) （若 \(j\) 为偶数），且 \(\Delta(j+1)\leqslant\log j/\log 2\) 对所有 \(j\) 成立）。借助 \(c\) ) 推出

\[
\mathrm{A} _ {2} (2 ^ {m}) \geqslant \frac {3}{2} r 2 ^ {m - 1} \geqslant \frac {3}{2} \frac {2 ^ {m} \log \log 2 ^ {m}}{2 \log 2}.
\]

由 \(c\) ) 与这个关系式断定： \(\mathrm{A}_2(n)\) 当 \(n\) 趋于 \(+\infty\) 时不等价于任何 (H) 函数。

\begin{enumerate}
\def\labelenumi{\arabic{enumi})}
\setcounter{enumi}{12}
\tightlist
\item
  设 \(f\) 是一个满足 \(1 \prec f(x) \prec x\) 的 (H) 函数。证明在泰勒公式
\end{enumerate}

\[
\begin{array}{l} f (x + f (x)) = f (x) + f (x) f ^ {\prime} (x) + \frac {1}{2 !} (f (x)) ^ {2} f ^ {\prime \prime} (x) + \dots + \\ \qquad + \frac {1}{n !} (f (x)) ^ {n} f ^ {(n)} (x) + \frac {1}{(n + 1) !} (f (x)) ^ {n + 1} f ^ {(n + 1)} (x + \theta f (x)) \end{array}
\]

（ \(0 < \theta < 1\) ，I, p.~42, exerc. 9）中，每一项相对于前一项都可忽略。若 f 相对于 x 的阶 \textless{} 1，则只要 n 充分大，这个和中的最后一项随 1/x 趋于 0。

\begin{enumerate}
\def\labelenumi{\arabic{enumi})}
\setcounter{enumi}{13}
\tightlist
\item
  由习题 13 推出：若 \(f\) 是一个使 \(\log f(e^x)\) 的阶 \(< 1\) （相对于 \(x\) ）的 (H) 函数，则函数 \(f(xf(x))\) 等价于形如 \(e^{g(x)}\) 的函数，其中 \(g\) 是 \(x\) 、 \(\log f(x)\) 以及后一函数的若干导数的有理函数。
\end{enumerate}

15) 设 \(f\) 是一个 (H) 函数，在 \(+\infty\) 的某个邻域上凸，且满足 \(f(x) \gg x\) 。对每个 \(\alpha > 0\) ，设 \(x_{\alpha}\) 是使 \(f'(x_{\alpha}) = (1 + \alpha)f'(x)\) 的点。

\begin{enumerate}
\def\labelenumi{\alph{enumi})}
\tightlist
\item
  若 \(f\) 的阶为 \(+\infty\) （相对于 \(x\) ），证明
\end{enumerate}

\[
x _ {\alpha} - x \sim \log (1 + \alpha) \frac {f (x)}{f ^ {\prime} (x)}
\]

（利用命题 2 (V, p.~251) 与 4 (V, p.~255)，把泰勒公式应用于 \(\log f'(x)\) ）。证明 \(\left(f(x_{\alpha}) - (x_{\alpha} - x)f'(x_{\alpha})\right)/f(x)\) 趋于 \((1 + \alpha)\left(1 - \log(1 + \alpha)\right)\) （当 \(x\) 趋于 \(+\infty\) 时）。

\begin{enumerate}
\def\labelenumi{\alph{enumi})}
\setcounter{enumi}{1}
\tightlist
\item
  若 \(f\) 的阶为 \(r > 1\) （相对于 \(x\) ），证明
\end{enumerate}

\[
x _ {\alpha} \sim (1 + \alpha) ^ {1 / (r - 1)} x
\]

（注意由 V, p.~255 的命题 4，若 \(f_{1}\) 的阶为 0（相对于 \(x\) ），则 \(f_{1}(kx) \sim f_{1}(x)\) 对每个常数 \(k\) 成立）。推出 \(\left(f(x_{\alpha}) - (x_{\alpha} - x)f'(x_{\alpha})\right) / f(x)\) 趋于

\[
r (1 + \alpha) - (r - 1) (1 + \alpha) ^ {r / (r - 1)}
\]

（当 \(x\) 趋于 \(+\infty\) 时）。

\begin{enumerate}
\def\labelenumi{\alph{enumi})}
\setcounter{enumi}{2}
\tightlist
\item
  最后设 \(f\) 相对于 \(x\) 的阶为 1。令 \(f(x) = xf_1(x)\) ，证明
\end{enumerate}

\[
x _ {\alpha} - x \sim \log (1 + \alpha) \frac {f _ {1} (x)}{f _ {1} ^ {\prime} (x)}
\]

（考察 \(f_{1}\) 的反函数，其阶为 \(+\infty\) （相对于 \(x\) ））。推出 \(\left(f(x_{\alpha}) - (x_{\alpha} - x)f'(x_{\alpha})\right) / f(x)\) 趋于 \(-\infty\) （当 \(x\) 趋于 \(+\infty\) 时）。

类似地，设 \(x_{\alpha}'\) 是使 \(f'(x_{\alpha}') = (1 - \alpha)f'(x)\) 的点（ \(0 < \alpha < 1\) ）。叙述与前面类似的表达 \(x - x_{\alpha}'\) 的主部以及

\[
\left(f (x _ {\alpha} ^ {\prime}) + (x - x _ {\alpha} ^ {\prime}) f ^ {\prime} (x _ {\alpha} ^ {\prime})\right) / f (x)
\]

（当 \(x\) 趋于 \(+\infty\) 时）的极限的公式。

由这些结果断定： \(f\) 是 \(+\infty\) 的某个邻域上的正则凸函数 (V, p.~260, exerc. 5)。

\makeatletter
\gdef\@chapapp{\chaptername}
\makeatother
\renewcommand{\thechapter}{\chinese{chapter}}
\renewcommand{\thesection}{\S\arabic{section}}
\setcounter{chapter}{5}
